\section{Introduction}

Score-debiased kernel density estimation (SD-KDE)~\cite{epstein2025sdkde}
improves the statistical efficiency of standard kernel density estimators:
for sufficiently smooth densities, SD-KDE attains better bias and mean-squared
error rates than vanilla KDE while retaining a simple, nonparametric form.
In particular, SD-KDE achieves an asymptotic mean integrated squared error
(AMISE) of order $O(n^{-8/(d+8)})$, improving the convergence exponent
relative to the $O(n^{-4/(d+4)})$ rate of classical KDE tuned with
Silverman's rule of thumb.
These accuracy gains come with a significant computational drawback: the
empirical score used for debiasing introduces an additional $O(n^2)$ pass
over the data, so a naive implementation has roughly the same quadratic cost
as KDE itself but with a larger constant factor.

Given samples $x_i$ and a bandwidth $h$, a Gaussian KDE is
\[
  \hat p(x)
  =
  \frac{1}{n h}
  \sum_{i=1}^n
  \varphi\!\left(\frac{x - x_i}{h}\right),
\]
and SD-KDE forms debiased samples
$x_i^{\mathrm{SD}} = x_i + \tfrac{h^2}{2}\,\hat s(x_i)$ where
$\hat s$ is an estimate of the score.  In this work we always use an
empirical score computed from the KDE itself, i.e.\ no parametric model
for the underlying density is assumed.  Writing the KDE explicitly in terms
of the samples,
\[
  \hat p(x)
  =
  \frac{1}{n h}
  \sum_{i=1}^n
  \varphi\!\left(\frac{x - x_i}{h}\right),
\]
we estimate the score as
\[
  \hat s(x)
  =
  \frac{\partial_x \hat p(x)}{\hat p(x)}
  =
  \frac{\sum_{i=1}^n \bigl[-\tfrac{x - x_i}{h}\,\varphi\!\left(\tfrac{x - x_i}{h}\right)\bigr]}
       {h^2 \sum_{i=1}^n \varphi\!\left(\tfrac{x - x_i}{h}\right)}.
\]
